% ID: SEL-D0005-OA-E01
% ESTADO: cotejada. Segunda pasada visual completa y borrador compilado revisados.
% Fuente: Selectividad/sel_2019_matematicas_aplicadas/Reserva A/Exámenes Mat Aplicadas CCSS - A y B.pdf
% SHA-256: d28871f7f203cf1be9e479e9d642cdb6b1623704a8b76a2b57cfa2d12950f293
% Página: 1; opción A; ejercicio 1. Mantener ambos apartados en una unidad.

\begin{enumerate}[label=\alph*)]
  \item \textbf{(1 punto)} Se considera el recinto cuadrado de vértices
  $(1,0)$, $(0,1)$, $(-1,0)$ y $(0,-1)$.

  Indique en qué puntos del recinto se alcanzan el valor máximo de la función
  $F(x,y)=3x+2y+7$ y el valor mínimo de la función $G(x,y)=x+y+6$,
  calculando dichos valores.
  \item \textbf{(1.5 puntos)} Resuelva la ecuación matricial
  $(A-A^t)\cdot X=B$, siendo $A$ y $B$ las matrices
  \[
    A=\begin{pmatrix}5&-2\\3&1\end{pmatrix}
    \qquad B=\begin{pmatrix}-1&3\\2&-1\end{pmatrix}
  \]
\end{enumerate}
